\documentclass[10pt]{article}
\usepackage[margin=1in]{geometry}
\usepackage{amsmath,amssymb,booktabs,array,microtype}
\usepackage{xcolor}
\usepackage[colorlinks=true,linkcolor=blue,urlcolor=blue]{hyperref}
\definecolor{navy}{HTML}{17365D}
\newcommand{\R}{\mathbb{R}}
\newcommand{\code}[1]{\texttt{\detokenize{#1}}}
\setlength{\parindent}{0pt}
\setlength{\parskip}{5pt}
\title{\color{navy}Curriculum Learning for Time-Series Users\\[3pt]\large Problem formulation and proposed treatment}
\author{Thesis project}
\date{Standalone method note}
\begin{document}
\maketitle
\begin{abstract}
This note isolates the scientific formulation from implementation and cluster
operations. It is intended to be readable on its own and reusable when drafting
a paper introduction or method section.
\end{abstract}

\section{Forecasting task}
Let $x_u(1{:}T)$ be the series of user $u\in\mathcal U$. At query time $t$,
the model receives $X_{u,t}=x_u(t-L+1{:}t)\in\R^L$ and predicts
$Y_{u,t}=x_u(t+1{:}t+H)\in\R^H$. Training and evaluation splits are
chronological and the complete target horizon belongs to its target split.

\section{Causal user difficulty}
Difficulty is computed from accessible training windows only. Persistence uses
$\widehat Y_{u,t,h}=x_u(t)$ and gives the normalized window error
\[
 d_{u,t}=\frac{\operatorname{MSE}(Y_{u,t},\widehat Y_{u,t})}
 {(\operatorname{std}(X_{u,t})+\varepsilon)^2}.
\]
The default user score is $d_u=\operatorname{median}_t d_{u,t}$. Users are
sorted by $(d_u,u)$, using the stable source identity to resolve ties.

\section{Proposed curriculum}
Training has $K$ phases and a fixed total of $S$ optimizer steps. Phase $k$
activates a cumulative prefix $A_k\subseteq\mathcal U$ of the ordered users.
At step $s$ in phase $k$, the sampler draws
\[
  u_s\sim\operatorname{Uniform}(A_k),\qquad
  (X_{u_s,t_s},Y_{u_s,t_s})\sim\mathcal W_{u_s}^{\rm train}.
\]
The model, optimizer, scheduler, and global step persist across phase
boundaries. Only the user sampling distribution changes.

\section{Controls and interpretation}
Uniform sampling tests the value of any curriculum. Hard-to-easy and a seeded
random order test whether the chosen order matters. The exposure-matched
control removes temporal order while retaining each user's expected number of
updates:
\[
 E_u=\sum_{k=1}^K S_k\frac{\mathbf 1[u\in A_k]}{|A_k|},\qquad
 p_u=\frac{E_u}{\sum_v E_v}.
\]
A gain over uniform but not over exposure matching indicates reweighting; a
gain over both supports a temporal-order effect.

\section{Reported evidence}
All methods share backbone, initialization distribution, optimizer budget,
batch size, loss, and evaluation rows. Reports include global error,
equal-user error, worst-10-percent-user error, convergence by optimizer step,
and error by difficulty quantile.
\end{document}
